% Generated draft: requires mathematical and source review.
\documentclass[11pt]{article}
\usepackage[margin=25mm]{geometry}
\usepackage{fontspec}
\setmainfont{Latin Modern Roman}
\usepackage{amsmath,amssymb,mathtools,unicode-math}
\setmathfont{Latin Modern Math}
\usepackage{enumitem,needspace,xurl,hyperref}
\setlength{\parindent}{0pt}
\setlength{\parskip}{5pt}
\setlength{\emergencystretch}{3em}
\newcommand{\N}{\mathbb{N}}
\newcommand{\Z}{\mathbb{Z}}
\newcommand{\Q}{\mathbb{Q}}
\newcommand{\R}{\mathbb{R}}
\newcommand{\C}{\mathbb{C}}
\newcommand{\F}{\mathbb{F}}
\newcommand{\A}{\mathbb{A}}
\newcommand{\PP}{\mathbb{P}}
\newcommand{\E}{\mathbb{E}}
\newcommand{\eps}{\varepsilon}
\newcommand{\dd}{\,\mathrm{d}}
\begin{document}
\section*{UnsolvedMath ID: 30005077}
\textbf{Source problem code:} catalogue entry
\subsection{Definitions and mathematical statement}
The supplied record identifies the target as Conjecture 2.5 of the cited Oberwolfach report, also labeled there as [2, Conjecture 6.7]. Let \(M\) denote the multigraded module, \(d\in\Z^r\) the truncation degree, and \(G^\bullet\) the complex referred to in that conjecture. The reconstructed statement calls \(G^\bullet\) a Fourier--Mukai complex; the original excerpt refers to its construction in an omitted part (1). Neither that construction nor its full hypotheses are supplied.

For clarity, if the intended truncation uses the coordinatewise order on \(\Z^r\), its notation means
\[
M_{\geq d}=\bigoplus_{\substack{a\in\Z^r\\ a_j\geq d_j\ (1\leq j\leq r)}}M_a.
\]
This interpretation is conditional: the excerpt does not explicitly define the truncation order or specify the graded coefficient ring that makes this sum a submodule.

\textbf{Conjecture, with unresolved scope.} Under the hypotheses of the omitted construction, the associated complex \(G^\bullet\) is a minimal graded free resolution of \(M_{\geq d}\).

More explicitly, writing \(F_i=G^{-i}\) in a homological indexing convention, the assertion requires a graded augmentation for which
\[
\cdots\longrightarrow F_2\xrightarrow{\partial_2}F_1
\xrightarrow{\partial_1}F_0\xrightarrow{\varepsilon}
M_{\geq d}\longrightarrow0
\]
is exact, each \(F_i\) is graded free over the intended coefficient ring, and the resolution is minimal. In a setting with the usual homogeneous maximal ideal \(\mathfrak m\), minimality means
\[
\partial_i(F_i)\subseteq\mathfrak m F_{i-1}\qquad(i\geq1).
\]
The indexing above only expresses the meaning of a resolution; the actual indexing, augmentation, coefficient ring, and applicable minimality convention must be recovered from the source construction.

The excerpt describes \(G^\bullet\) as a virtual resolution of \(M\), referring to [1], and reports a spectral sequence argument giving
\[
\beta_{i,a}(M_{\geq d})=\beta_{i,a}(G^\bullet).
\]
Here the excerpt does not supply the definition of the complex's Betti numbers or their index ranges. It also reports that the complex is a resolution in all examples considered by the authors. These are supplied contextual claims, not independently verified facts or substitutes for the conjectured exactness and minimality. The omitted construction and surrounding theorem are needed to determine whether restrictions involving \(d\)-regularity, the geometry, or other hypotheses belong in the quantifiers. Consequently, this is a conditional statement-only draft, not a fully specified universal assertion about arbitrary multigraded modules.
\subsection{Short English statement}
The conjecture asks whether the particular complex \(G^\bullet\) constructed in the cited report is the minimal graded free resolution of the truncation \(M_{\geq d}\). The excerpt reports matching Betti numbers and successful examples, but the construction and its hypotheses are missing. The precise scope therefore requires source review.
\subsection{Sources}
\begin{itemize}
\item Local UnsolvedMath registry, supplied source record: permanent ID \(30005077\), exact source problem code \texttt{catalogue entry}, title \emph{Minimal Resolutions from Virtual Multigraded Resolutions}, in the Oberwolfach Reports Open Problems collection. The record gives proposed year 2022 and names no proposer.
\item Supplied primary-source citation: \emph{Toric Geometry (2023)}, Oberwolfach Reports, DOI \texttt{10.4171/owr/2022/17}, pp.\ 861--927. \url{https://doi.org/10.4171/owr/2022/17}. The supplied excerpt labels the target Conjecture 2.5 and cross-references [2, Conjecture 6.7]; the bibliographic identity of [2] is not supplied.
\item Supplied literature source: Bruce--Cranton Heller--Sayrafi. \url{https://doi.org/10.1016/j.aim.2026.111049}. No title or passage establishing its precise relationship to the conjecture is supplied.
\item Supplied literature source: Oberwolfach Report 2022/17. \url{https://doi.org/10.4171/owr/2022/17}.
\item Literature assessment dated exactly \texttt{2026-08-22}: the supplied curation labels the problem open and states that later work characterizes regularity and virtual resolutions without establishing this exact claim in full generality. This is supplied dated curation, not independently verified literature assessment.
\item The registry labels the statement \texttt{corrected\_verified} and says it was reconstructed and checked against the official report. That verification is reported provenance; no source visit or independent verification is claimed in this draft.
\end{itemize}
\end{document}
